\documentclass[11pt]{scrartcl}

\usepackage[stixtwo]{fontsetup}
\usepackage{amsmath}

\begin{document}
Páros $n$ esetén
\[
	a_n
	= \left( \frac{n + 1}{n + 2} \right)^{n + 3}
	= \frac{\left(1 + \frac{1}{n}\right)^n}{\left(1 + \frac{2}{n}\right)^n}
	\cdot \left( \frac{n + 1}{n + 2} \right)^3
	\to \frac{1}{e},
\]
páratlan $n$ esetén
\[
	a_n
	= \left( \frac{n - 1}{n + 2} \right)^{n + 3}
	= \frac{\left(1 - \frac{1}{n}\right)^n}{\left(1 + \frac{2}{n}\right)^n}
	\cdot \left( \frac{n - 1}{n + 2} \right)^3
	\to \frac{1}{e^3},
\]
így a sorozat torlódási pontjainak halmaza
$\left\{ \frac{1}{e}, \frac{1}{e^3} \right\}$,
$\limsup a_n = \frac{1}{e}$,
$\liminf a_n = \frac{1}{e^3}$, és mivel
$\liminf a_n \ne \limsup a_n$, így a sorozat divergens.
\end{document}
